\section{The problem, and why it is hard}
\label{sec:intro}

\subsection{What we are trying to decide}

Given two SQL queries $Q_1$ and $Q_2$ over the same schema, decide whether they return
the same answer \emph{on every possible database instance}. Not on your data --- on all data.

This matters because every relational engine contains a rule-based rewriter. Calcite,
Spark SQL, CockroachDB and PostgreSQL all transform your query into a different one before
executing it. If a rewrite rule is wrong, the engine silently returns wrong answers.
Historically these bugs have been found by users, years later~\cite{udp}.

The general problem is undecidable. For queries with negation and arbitrary predicates you
can encode first-order validity, so no algorithm decides all cases. Everything in this
document is therefore a \emph{sound but incomplete} procedure: when it says
``equivalent'' it is right; when it says ``unknown'' the queries may still be equivalent.

\keybox{The three answers a prover may give}{%
\textbf{EQUIVALENT} --- proved, trust it.\quad
\textbf{NOT EQUIVALENT} --- a counterexample database exists.\quad
\textbf{UNKNOWN} --- the procedure ran out of power.
\\[3pt]
All three systems here answer EQUIVALENT or UNKNOWN. None of them is a
counterexample finder. Refutation is a different problem, usually attacked
by bounded model finding.}

\subsection{Bag semantics, and why it changes everything}

SQL tables are \emph{bags} (multisets): duplicates are real and observable.
\lstinline|SELECT x FROM R| does not deduplicate. Two queries that agree as sets may
disagree as bags:

\begin{lstlisting}[style=sql]
SELECT x FROM R                 -- may return 1, 1, 2
SELECT DISTINCT x FROM R        -- returns    1, 2
\end{lstlisting}

The classical theory of query equivalence (containment mappings, the chase, homomorphism
checking) is set-based. Under bags the picture is harder: bag \emph{containment} of
conjunctive queries is a long-standing open problem, while bag \emph{equivalence} of unions
of conjunctive queries is decidable~\cite{cohen}. That decidability result is the
theoretical anchor for all three systems.

\subsection{The common architecture}

All three systems share a pipeline:

\begin{center}
\begin{tikzpicture}[node distance=6mm,
  box/.style={draw=cGrey,fill=cGreyBg,rounded corners=3pt,inner sep=5pt,
              font=\footnotesize\sffamily,minimum height=9mm,align=center},
  hi/.style ={draw=cBlue,fill=cBlueBg,rounded corners=3pt,inner sep=5pt,
              font=\footnotesize\sffamily,minimum height=9mm,align=center},
  ar/.style ={-{Stealth[length=2.4mm]},cGrey,thick}]
\node[box] (sql) {SQL text};
\node[box,right=of sql] (plan) {relational\\plan};
\node[hi,right=of plan]  (alg) {algebraic form\\(U- or Q-expression)};
\node[hi,right=of alg]  (norm) {normal form};
\node[box,right=of norm] (dec) {decide};
\draw[ar] (sql)--(plan); \draw[ar] (plan)--(alg);
\draw[ar] (alg)--(norm); \draw[ar] (norm)--(dec);
\end{tikzpicture}
\end{center}

They differ in exactly two places, and it is worth stating this up front because everything
else follows from it:

\begin{enumerate}
\item \textbf{The algebraic form.} \udp{} and \sqlsolver{} use \emph{U-expressions},
      defined \emph{axiomatically} over an abstract structure called a U-semiring.
      \QED{} uses \emph{Q-expressions}, which are the same syntax pinned to a
      \emph{concrete} semantic domain $\Nbar = \mathbb{N}\cup\{\infty\}$.
\item \textbf{The decision step.} \udp{} compares normal forms \emph{syntactically}.
      \QED{} encodes them into first-order logic and calls an SMT solver.
      \sqlsolver{} encodes them into \liastar{}, a theory built for unbounded sums, and
      reduces that to plain \lia{} for an SMT solver.
\end{enumerate}

\begin{table}[h]
\centering\small
\begin{tabularx}{\linewidth}{@{}l l l X@{}}
\toprule
 & \textbf{algebra} & \textbf{decision step} & \textbf{one-line character} \\
\midrule
\udp{} 2018      & U-expression & syntactic isomorphism & the foundation; proves things by axiom-chasing \\
\QED{} 2024      & Q-expression & SMT on FOL encoding   & better \emph{modelling}: NULLs, constraints, unknown operators \\
\sqlsolver{} 2024& U-expression & \liastar{} $\to$ \lia{} $\to$ SMT & better \emph{arithmetic}: unbounded sums, real aggregates \\
\bottomrule
\end{tabularx}
\caption{The three systems at a glance. \spes{}~\cite{spes}, a fourth system, uses a
\emph{tree-structured} algebraic representation of relational operators rather than a
U-expression, normalises it with hand-written rewrite rules, and calls an SMT solver only to
compare predicates; we use it as a baseline throughout.}
\end{table}

\subsection{A warning about the published numbers}

Every paper reports how many query pairs it proves on ``the Calcite test suite''. These
numbers are \emph{not} comparable across papers, because each group extracted its own pairs
from a different Calcite version with a different filter. The clearest evidence: \QED{} and
\sqlsolver{} both report \spes{} on ``232 Calcite pairs'' and get \textbf{95} and
\textbf{99} respectively. Treat cross-paper percentages as indicative only. We return to
this in Section~\ref{sec:compare}.
